\begin{derivation}[全入射振幅与物理散射通道]
全入射约定把每一条外腿的四动量都定义成指向相互作用区。设全入射核的四个参数为 $k_i$，则正文\eqnrefs{eq:all-incoming-amplitude-dp9} 要求
\[
k_1+k_2+k_3+k_4=0.
\]
对物理过程 $1+2\to3+4$，实验四动量满足
\[
p_1+p_2=p_3+p_4.
\]
前两条腿本来就流入，所以
\[
k_1=p_1,\qquad k_2=p_2.
\]
末态 3、4 的物理箭头流出相互作用区；若仍要求 $k_i$ 全部向内，必须反号：
\[
k_3=-p_3,\qquad k_4=-p_4.
\]
于是全入射的动量守恒自动变成
\begin{align*}
k_1+k_2+k_3+k_4
&=p_1+p_2-p_3-p_4\\
&=0,
\end{align*}
与物理守恒完全一致。将四个 $k_i$ 代回同一个截肢四点核，便得到正文\eqnrefs{eq:physical-from-all-incoming-dp9}：
\[
\mathcal M_{12\to34}
=\mathcal A(p_1,p_2,-p_3,-p_4).
\]
\end{derivation}

\begin{derivation}[低速散射振幅与势能的匹配]
低速弹性区有 $|\vb{p}_f|=|\vb{p}_i|$。定义非相对论振幅为阈值极限
\begin{equation*}
\mathcal M_{\rm rel}^{\rm NR}(\vb{q})
\equiv\lim_{|\vb{p}_i|/(m_{A,B}c)\to0}\mathcal M_{\rm rel},
\qquad
\lim_{|\vb{p}_i|\to0}\sqrt{s}=(m_A+m_B)c.
\end{equation*}
因此对截面主式取同一极限得到
\begin{equation*}
 \frac{\dd\sigma}{\dd\Omega}
 =\frac{\hbar^2}{64\pi^2(m_A+m_B)^2c^2}
 |\mathcal M_{\rm rel}^{\rm NR}(\vb{q})|^2.
\end{equation*}
另一方面，对\eqnrefs{eq:sch-lippmann-r13} 作一次 Born 迭代，采用同一出射波相位约定可写
\begin{equation*}
 f_B(\vb{q})
 =-\frac{m_{\rm red}}{2\pi\hbar^2}
 \widetilde V_{\rm NR}(\vb{q}),
 \qquad
 \frac{\dd\sigma}{\dd\Omega}=|f_B|^2,
\end{equation*}
其中 $m_{\rm red}=m_Am_B/(m_A+m_B)$。与相对论截面对应的散射振幅归一化为
\begin{equation*}
 f_{\rm rel}(\vb{q})
 =\frac{\hbar}{8\pi(m_A+m_B)c}
 \mathcal M_{\rm rel}^{\rm NR}(\vb{q}),
\end{equation*}
其相位由前面固定的 $S$ 矩阵约定与出射 Green 核共同确定。令 $f_B=f_{\rm rel}$，得到
\begin{align*}
 \widetilde V_{\rm NR}(\vb{q})
 &=-\frac{\hbar^3}{4m_{\rm red}(m_A+m_B)c}
 \mathcal M_{\rm rel}^{\rm NR}(\vb{q})\\
 &=-\frac{\hbar^3}{4m_Am_Bc}
 \mathcal M_{\rm rel}^{\rm NR}(\vb{q}),
\end{align*}
即\eqnrefs{eq:relativistic-to-nr-potential-matching-dp75}。
\end{derivation}

\begin{derivation}[质心系中的二体相空间]
正文\eqnrefs{eq:phase2-si-start-r9} 仍含两个末态三动量积分和一个四维守恒 $\delta$ 函数。在质心系中，三维 $\delta$ 函数先固定 $\vb{p}_4=-\vb{p}_3$，剩下的能量 $\delta$ 函数再固定径向动量模；最后只留下末态方向 $\dd\Omega$，从而得到正文\eqnrefs{eq:phase2-final-si-r9}。

采用能量四矢量
\begin{equation*}
 \mathcal P^\mu\equiv cp^\mu=(E,c\vb{p}),
 \end{equation*}
它只是物理四动量乘以常数 $c$。定义标准二体相空间因子
\begin{equation*}
 \dd\mathfrak P_2
 =(2\pi)^4\delta^{(4)}(\mathcal P_1+\mathcal P_2-\mathcal P_3-\mathcal P_4)
 \prod_{a=3,4}\frac{c^3\dd^3 p_a}{(2\pi)^3\,2E_a}.
 \end{equation*}
每个 $c^3\dd^3 p$ 都是 $\dd^3(c\vb{p})$，所以 \eqnrefs{eq:phase2-si-start-r9} 完全由能量量纲的四动量组成。在质心系 $\mathcal P_1+\mathcal P_2=(c\sqrt s,\mathbf0)$。三维 $\delta$ 先令 $\vb{p}_4=-\vb{p}_3$，于是
\begin{equation*}
 \dd\mathfrak P_2
 =\frac{c^4p^2\dd p\,\dd\Omega}{(2\pi)^2 4E_3E_4}
 \delta(c\sqrt s-E_3-E_4).
\end{equation*}
令
\begin{equation*}
 f(p)=c\sqrt s-E_3(p)-E_4(p).
\end{equation*}
由 $\dd E_a/\dd p=c^2p/E_a$，
\begin{equation*}
 \left|f'(p)\right|
 =c^2p\left(\frac1{E_3}+\frac1{E_4}\right)
 =\frac{c^3p\sqrt s}{E_3E_4}.
\end{equation*}
所以
\begin{equation*}
 \delta[f(p)]
 =\frac{E_3E_4}{c^3p_f\sqrt s}\delta(p-p_f).
\end{equation*}
代回并积分 $p$，得到
\begin{equation*}
 {
 \dd\mathfrak P_2
 =\frac1{16\pi^2}\frac{|\vb{p}_f|}{\sqrt s}\,\dd\Omega.}
 \end{equation*}
右端无量纲，因为 $|\vb{p}_f|$ 与 $\sqrt s$ 都具有动量量纲。
\end{derivation}

\begin{derivation}[不变入射通量与相遇速率]
对两条入射轨迹，普通同一参考系中的相对相遇速率不能简单写成任意参考系都适用的 $|\vb{v}_1-\vb{v}_2|$。由四动量构造 Lorentz 标量

\begin{equation*}
 \mathcal I_{12}
 \equiv\sqrt{(p_1\!\cdot p_2)^2-m_1^2m_2^2c^4}.
 \end{equation*}
它具有动量平方量纲。定义 M\o ller 相对速率
\begin{equation*}
 v_{\rm M}
 \equiv
 \frac{c^3\mathcal I_{12}}{E_1E_2}.
 \end{equation*}
检查质心系时不跳过中间平方。写
\begin{equation*}
 p_1^\mu=\left(\frac{E_1}{c},\vb{p}_i\right),
 \qquad
 p_2^\mu=\left(\frac{E_2}{c},-\vb{p}_i\right),
\end{equation*}
于是
\begin{equation*}
 p_1\!\cdot p_2=\frac{E_1E_2}{c^2}+|\vb{p}_i|^2.
\end{equation*}
两条质量壳又分别给
\begin{equation*}
 m_a^2c^2=\frac{E_a^2}{c^2}-|\vb{p}_i|^2,
 \qquad a=1,2.
\end{equation*}
因此 \eqnrefs{eq:invariant-flux-scalar-r9} 根号中的量逐项化为
\begin{align*}
 \mathcal I_{12}^2
 &=\left(\frac{E_1E_2}{c^2}+p_i^2\right)^2
 -\left(\frac{E_1^2}{c^2}-p_i^2\right)
  \left(\frac{E_2^2}{c^2}-p_i^2\right),\\
 &=\frac{E_1^2E_2^2}{c^4}
 +\frac{2E_1E_2p_i^2}{c^2}+p_i^4
 -\frac{E_1^2E_2^2}{c^4}
 +\frac{(E_1^2+E_2^2)p_i^2}{c^2}-p_i^4,\\
 &=\frac{p_i^2}{c^2}(E_1+E_2)^2,
\end{align*}
其中 $p_i\equiv|\vb{p}_i|\ge0$。取正根得到
\begin{equation*}
 {
 \mathcal I_{12}
 =\frac{|\vb{p}_i|(E_1+E_2)}{c}
 =|\vb{p}_i|\sqrt s.}
 \end{equation*}
最后一步用了 $E_1+E_2=c\sqrt s$。于是
\begin{align*}
 v_{\rm M}
 &=|\vb{p}_i|c^2
 \left(\frac1{E_1}+\frac1{E_2}\right),\\
 &=|\vb{v}_1-\vb{v}_2|,
\end{align*}
因为 $\vb{v}_a=c^2\vb{p}_a/E_a$。所以 \eqnrefs{eq:moller-velocity-si-r9} 的协变表达确实回到熟悉的相对速度。
\end{derivation}

\begin{derivation}[盒归一化怎样给出二体散射截面]

取空间周期盒体积 $V$ 与观测时间 $T$。空间连续态满足
\begin{equation*}
 (2\pi\hbar)^3\delta^{(3)}(\mathbf0)=V,
 \qquad
 \sum_{\vb{p}}\longrightarrow
 V\int\frac{\dd^3 p}{(2\pi\hbar)^3}.
\end{equation*}
时间方向令 $\delta E\equiv E_f-E_i$，则
\begin{equation*}
 \int_{-T/2}^{T/2}\dd t\,\ee^{\ii\delta E t/\hbar}
 =T\,\operatorname{sinc}\!\left(\frac{\delta E T}{2\hbar}\right),
\end{equation*}
并由正文\eqnrefs{eq:delta-square-time-r15} 有
\begin{equation*}
 \lim_{T\to\infty}\frac1T
 \left|\int_{-T/2}^{T/2}\dd t\,\ee^{\ii\delta E t/\hbar}\right|^2
 =2\pi\hbar\delta(\delta E).
\end{equation*}
因此四维动量守恒 $\delta$ 函数平方在概率除以 $T$ 后留下一个能量守恒 $\delta$，其空间零点因子给出一个 $V$。结合两条初态的协变归一化，单位盒中单对入射粒子的二体末态跃迁率为正文\eqnrefs{eq:box-transition-rate-r15}，即
\begin{equation*}
 \dd\Gamma_{1+2\to2}
 =\frac{\hbar^2c^3}{4VE_1E_2}
 \overline{|\mathcal M|^2}\,\dd\mathfrak P_2.
\end{equation*}

单个归一化入射粒子在盒中的数密度是 $1/V$。正文\eqnrefs{eq:moller-velocity-si-r9} 给出相对速率，所以入射通量密度为
\begin{equation*}
 j_{12}=\frac{v_{\rm M}}{V}
 =\frac{c^3\mathcal I_{12}}{VE_1E_2}.
\end{equation*}
按正文中“跃迁率除以入射通量密度”的截面关系：
\begin{align*}
 \dd\sigma
 &=\frac{\dd\Gamma_{1+2\to2}}{j_{12}}\\
 &=\frac{\hbar^2}{4\mathcal I_{12}}
 \overline{|\mathcal M|^2}\,\dd\mathfrak P_2.
\end{align*}
$V$、$T$ 与参考系依赖的 $E_1E_2$ 在这一步全部消去。对 $2\to2$ 过程代入正文\eqnrefs{eq:phase2-final-si-r9}，并在质心系使用
$\mathcal I_{12}=|\vb{p}_i|\sqrt s$，得到
\begin{align*}
 \frac{\dd\sigma}{\dd\Omega}
 &=\frac{\hbar^2}{4|\vb{p}_i|\sqrt s}
 \overline{|\mathcal M|^2}
 \frac1{16\pi^2}\frac{|\vb{p}_f|}{\sqrt s}\\
 &=\frac{\hbar^2}{64\pi^2s}
 \frac{|\vb{p}_f|}{|\vb{p}_i|}
 \overline{|\mathcal M|^2},
\end{align*}
即正文\eqnrefs{eq:2to2-cross-section-master-si-r9}。
\end{derivation}

\begin{derivation}[单粒子衰变率的归一化]
协变态归一化的单粒子初态在盒归一化后带因子 $(2P^0V)^{-1/2}$； 末态每条外腿给对应的质量壳测度。 LSZ 矩阵元中的四维 $\delta$ 在平方后，空间部分利用

\begin{equation*}
 (2\pi\hbar)^3\delta^{(3)}(\mathbf0)=V
\end{equation*}
消去母粒子的盒体积，而时间部分利用 \eqnrefs{eq:delta-square-time-r15} 把 $|\delta(E_f-E_i)|^2/T$ 化为单个能量 $\delta$。于是概率除以 $T$ 后只剩一个母粒子归一化因子 $1/(2P^0)$ 与 $n$ 体末态相空间。

为避免混淆能量量纲和动量量纲，先定义能量归一化截肢振幅
\begin{equation*}
 \mathcal A_E\equiv c\,\mathcal M.
\end{equation*}
若 $\mathcal M$ 在某个过程具有一个动量量纲，则 $\mathcal A_E$ 具有能量量纲。有限盒推导给
\begin{equation*}
 \dd\Gamma_{1\to n}
 =\frac{1}{2E_P\hbar}\,
 \overline{|\mathcal A_E|^2}\,\dd\mathfrak P_n,
 \qquad E_P=Mc^2.
 \end{equation*}
把 $\mathcal A_E=c\mathcal M$ 与 $E_P=Mc^2$ 代入，两个 $c^2$ 消去：
\begin{align*}
 \dd\Gamma_{1\to n}
 &=\frac{c^2}{2Mc^2\hbar}
 \overline{|\mathcal M|^2}\,\dd\mathfrak P_n,\\
 &={
 \frac{1}{2M\hbar}
 \overline{|\mathcal M|^2}\,\dd\mathfrak P_n.}
 \end{align*}
这里 $\dd\mathfrak P_n$ 使用能量四矢量 $\mathcal P^\mu=cp^\mu$ 的定义。以无量纲 Yukawa 耦合的费米子对衰变为例，振幅含有双线性因子 $\bar u v$；在上述外腿归一化下该双线性具有动量量纲，因此
\begin{equation*}
 \left[\frac{|\mathcal M|^2}{M\hbar}\right]
 =\mathrm{s^{-1}},
\end{equation*}
与衰变率的 SI 单位一致。这样 \eqnrefs{eq:master-decay-rate-si-r9} 与散射截面主式共享同一套有限盒、有限时间和外腿归一化，而不是独立记忆的另一条规则。
\end{derivation}

\begin{derivation}[三体相空间与 Dalitz 变量]
在母粒子静止系中，对正文\eqnrefs{eq:phase3-recursive-dp52} 的两个二体因子分别使用\eqnrefs{eq:phase2-final-si-r9}：

\begin{equation*}
\dd\mathfrak P_3
=\frac{\dd\mathcal S_{12}}{2\pi}
\frac{1}{16\pi^2}\frac{|\vb{p}_3|}{Mc}\dd\Omega_3
\frac{1}{16\pi^2}\frac{|\vb{p}_2^*|}{\sqrt{\mathcal S_{12}}/c}\dd\Omega_2^*.
\end{equation*}
若母粒子无自旋或已做自旋平均，整体方向积分给 $\int\dd\Omega_3=4\pi$。在 $q_{12}$ 静止系中以粒子 3 为极轴，方位角再给 $2\pi$，于是
\begin{equation*}
\dd\mathfrak P_3
=\frac{\dd\mathcal S_{12}\,\dd\cos\theta^*}{64\pi^3}
\frac{|\vb{p}_3|}{Mc}
\frac{|\vb{p}_2^*|}{\sqrt{\mathcal S_{12}}/c}.
\end{equation*}
正文\eqnrefs{eq:dalitz-angular-map-dp68} 给出
\begin{equation*}
\left|\frac{\partial\mathcal S_{23}}{\partial\cos\theta^*}\right|
=2c^2|\vb{p}_2^*||\vb{p}_3^*|.
\end{equation*}
因此把 $\cos\theta^*$ 换成 $\mathcal S_{23}$ 时，两个动量模恰好与 Jacobian 相消；再用 $\sqrt{\mathcal S_{12}}$ 和母粒子质量整理剩余因子，得到
\begin{equation*}
\dd\mathfrak P_3
=\frac{\dd\mathcal S_{12}\,\dd\mathcal S_{23}}{128\pi^3M^2c^4},
\end{equation*}
即正文\eqnrefs{eq:phase3-dalitz-dp52}。
\end{derivation}

\begin{derivation}[二体阈值附近的相空间因子]

正文\eqnrefs{eq:bubble-imag-phase-space-dp8} 表明虚部只有在两个内部粒子同时满足正能质量壳和能量守恒时才存在。对等质量情形进入质心系，能量 $\delta$ 函数固定唯一的径向动量 $k_*$；解出这个根并把它写成相对论速度因子，就得到正文\eqnrefs{eq:two-body-threshold-beta-dp8}。

在质心系 $\vb{P}=0$ 且 $m_1=m_2=m$ 时，
\begin{equation*}
 E_1=E_2=E_k=\sqrt{m^2c^4+c^2k^2},
 \qquad
 \sqrt{s}\,c=2E_k.
\end{equation*}
平方后有
\begin{align*}
 sc^2
 &=4m^2c^4+4c^2k_*^2,\\
 k_*^2
 &=\frac{s-4m^2c^2}{4}.
\end{align*}
因此只有 $s\ge4m^2c^2$ 时 $k_*$ 才为实数。再用 $E_{\rm cm}=\sqrt{s}\,c$，
\begin{align*}
 \beta_*(s)
 &=\frac{ck_*}{E_{\rm cm}/2}
 =\frac{2ck_*}{\sqrt{s}\,c}\\
 &=\sqrt{1-\frac{4m^2c^2}{s}},
\end{align*}
即正文\eqnrefs{eq:two-body-threshold-beta-dp8}。为把阈值因子的整体归一化也固定下来，定义正文\eqnrefs{eq:bubble-imag-phase-space-dp8} 中的径向相空间核
\begin{equation*}
 \rho_2(s)
 \equiv
 \int\frac{\dd^3\vb{k}}{(2\pi\hbar)^3}
 \frac{c^2}{4E_k^2}
 \delta\!\left(\sqrt{s}\,c-2E_k\right).
\end{equation*}
利用
\begin{equation*}
 \left.\frac{\dd(2E_k)}{\dd k}\right|_{k_*}
 =\frac{2c^2k_*}{E_*},
 \qquad E_*\equiv E_{k_*}=\frac{\sqrt{s}\,c}{2},
\end{equation*}
径向 $\delta$ 积分严格给出
\begin{align*}
 \rho_2(s)
 &=\frac{4\pi k_*^2}{(2\pi\hbar)^3}
 \frac{c^2}{4E_*^2}
 \frac{E_*}{2c^2k_*}\\
 &=\frac{k_*}{16\pi^2\hbar^3E_*}
 =\frac{\beta_*(s)}{16\pi^2\hbar^3c}.
\end{align*}
因此连续谱在阈值处按确定的线性因子 $\beta_*(s)$ 从零打开。
\end{derivation}
